\PassOptionsToPackage{obeyspaces}{url}
\documentclass[UTF8,fontset=windows,11pt,a4paper]{ctexart}
\usepackage[margin=24mm]{geometry}
\usepackage{amsmath,amssymb,booktabs,tabularx,graphicx,float}
\usepackage[hidelinks]{hyperref}
\usepackage{url}
\setmainfont{Times New Roman}
\setmonofont{Consolas}
\setlength{\parskip}{0.35em}
\setlength{\parindent}{2em}
\setlength{\emergencystretch}{3em}
\hypersetup{pdftitle={数据挖掘作业逐题讲解},pdfauthor={}}
\title{数据挖掘作业逐题讲解}
\author{}
\date{2026年11月26日}
\begin{document}
\maketitle
\tableofcontents
\newpage

\section*{阅读说明}
\addcontentsline{toc}{section}{阅读说明}
这份讲解与第一题正式报告分开：正式报告面向提交和查证，本稿解释为什么这样建模、每一步输入输出是什么，以及结果能够支持什么结论。第一题依据本次实际运行；第二、三、四题和拓展一依据队友现有代码与保存结果，本次没有重新训练这些题。拓展二只说明题目设计路线，未实现系统。

\begin{center}\small
\begin{tabularx}{\linewidth}{lXXX}\toprule
任务 & 输入与输出 & 可用监督信息 & 核心判断\\\midrule
第一题 & 图像到六个簇 & 课程标签只作最终评价 & 哪些图像相似\\
第二题 & 图像到异常分数 & 训练正常、异常标签均可用 & 图像是否有缺陷\\
第三题 & 历史窗口到下一时刻数值 & 历史目标值用于回归训练 & 下一时刻是多少\\
第四题 & 六维记录到异常分数 & 拟合数据只有正常样本 & 是否偏离正常规律\\
拓展一 & 两种模态到异常分数 & 各模态分别规定 & 能否共用正常参考流程\\\bottomrule
\end{tabularx}\end{center}

理解这些题时始终分开三件事：\textbf{表示}决定保留哪些信息，\textbf{模型}决定如何利用信息，\textbf{评价协议}决定结果是否可信。模型更复杂不自动意味着效果更好；有较高分数也不自动说明对新数据有效。

\section{第一题：从图片到六个无标签分组}
\subsection{题目到底要求什么}
数据包有600张图，属于瓶子、电缆、皮革、药片、瓷砖、晶体管六类，每类100张。真实类别保存在JSON中，但聚类程序在完成所有分组并保存文件之后才读取它，用来检查结果。聚类输出是簇编号，例如0到5；编号本身没有“瓶子”等语义。

题面给出六类，因此可以直接设置$K=6$。虽然每类100张是已知事实，本实现没有强迫每个簇恰好100张；分组大小由特征距离和算法产生。全部600张参与无标签聚类与PCA拟合，没有虚构训练／测试划分。这回答给定图片如何分组，并未测试未知新图像。

评分中的形式化、特征、算法、评价分别占5、5、10、5分。完整回答需要解释选择理由、给出对照与结果分析，并提供能够重跑的代码。

\subsection{为什么先把图片变成向量}
原始像素很多，而且位置稍有变化就会导致大量像素不同。皮革和瓷砖主要由纹理区分，瓶子、电缆和晶体管还有整体形状差异。我们先用特征向量描述这些信息，再比较向量距离。一个图像对应一行，整个数据集变成$600\times d$的矩阵。

\textbf{传统基线。}把图像缩放到$224\times224$。在HSV颜色空间按$8\times4\times4$分箱，统计128维颜色频率；再从灰度图计算uniform LBP，使用8个邻居、半径1，得到10维局部纹理直方图。LBP比较中心像素与邻居亮暗关系，用这些局部模式出现的频率描述纹理。颜色频率不保留具体空间位置，因此易于比较，但可能把颜色相近、形状不同的物体混在一起。

两个直方图先各自归一化为总和1，再逐项开方，最后等权拼接并除以$\sqrt2$，得到138维单位向量。这使两个特征块都具有相同总权重，并通过Hellinger形式调整频率差异的影响。128维颜色块不会仅因为维度更高就占据全部权重。

\textbf{深度特征。}使用ImageNet预训练的ResNet18，去掉最后的分类层，保留全局平均池化输出的512维向量。网络参数冻结，不在这600张图上更新。它提供的是外部预训练的视觉表示，而不是直接调用其原来的1000类分类结果。预处理使用该权重的官方约定：短边缩放到256、中心裁剪224，再按ImageNet通道均值与标准差处理。

这种方式准确说是“外部有监督预训练特征上的无监督聚类”。传统基线没有这种外部标签来源。两条管线的预处理也不同，因此对照评价的是完整特征方案，不能把全部差异归因于某一个网络层。

\subsection{归一化与PCA分别在做什么}
对深度向量$z$执行$z/\|z\|_2$，使距离主要反映特征方向。对于两个单位向量，$\|u-v\|_2^2=2-2u^Tv$，因此两者越相似，欧氏距离越小。这里的普通K-means仍然更新普通质心，并没有把质心强制成单位向量，所以不称为球面K-means。

PCA对照先在所有无标签深度向量上找到方差较大的方向，保留64维。它可能压缩冗余，也可能舍弃恰好有助于区分类别的信息，所以不能预设降维一定提高成绩。PCA会中心化并投影，投影结果不再保持单位长度，因此降维后再做一次L2归一化。不进行白化。主方法使用原始512维，PCA是单独对照，不根据真实标签决定是否采用。

\subsection{K-means如何得到分组}
目标函数为
\[
\min_{\{c_i\},\{\mu_k\}}\sum_{i=1}^{600}\|z_i-\mu_{c_i}\|_2^2.
\]
算法交替进行两步：把每张图分给最近的中心；用该簇所有向量的均值更新中心。重复直到收敛或达到迭代上限。K-means++用于初始化中心，降低初始中心挤在一起的风险；一次运行使用20次初始化，并按不含真实标签的簇内平方距离目标保留最优者。

三种K-means表示分别使用11、23、37、53、71五个预先固定的种子。主方法和主结果种子11在运行前确定，五个种子的结果全部报告。另一组对照使用同一512维深度特征上的Ward凝聚层次聚类：从各点独立成簇开始，每次合并使簇内平方和增量最小的两个簇，直到剩六簇。它不依赖K-means的随机初始化，但仍受特征空间和距离定义影响。

\subsection{如何判断分组是否正确}
簇编号可以任意交换。例如所有瓶子被标为簇3，只要它们聚在一起就不是错误。ACC先构造“簇×真实类别”的计数表，再求一对一匹配，使对角线总数最大；最后除以600。不能直接比较原始编号，也不能让多个簇随意映射到同一类来抬高准确率。

ARI从样本对出发，比较“同类／不同类”关系，并校正随机一致性；NMI从信息量出发，比较两种分组共享多少信息。二者不需要重排簇编号。跨随机种子的稳定性也可用ARI，不过比较的是两次预测分组，而非预测与真实标签。

轮廓系数比较一个点与本簇、最近其他簇的距离，CH和DB分别从簇间分离与簇内散布衡量结构。这些指标不用标签，但不等于语义正确率；不同表示空间的数值不能直接当作统一排名。二维PCA图只是投影可视化，实际聚类在138、512或64维中进行。

\subsection{本次实际结果与解释}
\InputIfFileExists{task1-actual-results.tex}{}{}

\subsection{怎样复现、查错和讲清楚}
先读\path{Problem 1/README.md}，再看\path{features.py}的表示、\path{run.py}的固定协议和聚类、\path{verify.py}的独立评价。主要证据依次为\path{protocol.json}、\path{clusters.npz}、\path{predictions.csv}、\path{all_metrics.csv}。在读取标签之前保存分组和文件哈希，有助于检查单次执行顺序；哈希本身不能证明整个研究历史，也不能证明新数据上的泛化。

讲解这题可以按五句话组织：数据是什么；如何提特征；为什么用六簇K-means；如何通过对照区分特征和算法作用；用哪些指标与实例支持结论。答辩时应能解释“冻结网络不等于整个过程没有监督来源”“聚类ACC为什么需要匹配”“每簇100张是否人为强制”“为什么没有训练测试划分”“为什么高分不能直接外推”。
\section{第二题：图像异常检测}

\subsection{问题、数据和输出}

给定一张榛子或拉链图像，判断它是否存在异常。产品类别事先已知，两个产品分别建模，不需要先判断图片属于哪种产品。每类训练集都有 200 张正常图、50 张异常图；榛子测试集有 40 张正常、15 张异常，拉链有 32 张正常、15 张异常。训练图共 500 张，测试图共 102 张。

输出首先是连续异常分数，再与该产品的阈值比较，分数大于等于阈值判为异常。因此，这题同时涉及“把异常排在前面”和“确定实际判断边界”。训练集中存在异常标签，现有整体方案属于使用监督信息的异常分类；只有其中的正常记忆库分支只学习正常图像。

依据：\path{Problem 2/course-report/sec/task2-experiments.tex:3}；\path{Problem 2/course-report/sec/task2-method.tex:24}。

\subsection{从图片变成可计算的特征}

原始图像为 1024×1024，代码先转 RGB，再把完整画面缩放至 384×384，并使用 ImageNet 的通道均值、标准差处理像素。随后送入冻结的预训练 Wide ResNet50-2。“冻结”表示不在这 500 张训练图上更新卷积网络参数，而是借用网络已有的视觉表示，再训练较小的判别模块。

同一次网络前向计算产生两种表示：

\par\noindent\textbullet\quad \textbf{全局表示}：对 layer2、layer3、layer4 分别计算各通道的空间均值和最大值，再拼接，得到每张图 7168 维向量。均值描述整体外观，最大值保留突出响应。\par
\par\noindent\textbullet\quad \textbf{局部表示}：保留 layer2、layer3 的空间结构，平滑并对齐网格，拼接后用固定随机矩阵由 1536 维投影到 256 维，得到 48×48 个局部向量。每张图对应形状为 2304×256 的描述子集合。\par

全局向量可以把整张图交给普通分类器；局部向量保留了缺陷可能出现的位置。两者解决的是信息粒度不同的问题。随机投影只负责降低存储和距离计算成本，不是通过标签学习出来的新网络。

依据：\path{Problem 2/course-report/sec/task2-method.tex:27}、同文件第29行、同文件第50行。

\subsection{两个检测分支如何工作}

\textbf{全局分支使用 RBF 核 SVM。}先在当前拟合数据上估计每个特征的均值、标准差，再对特征标准化；SVM 利用正常和异常训练标签学习边界。RBF 核为
\[
K(u,v)=\exp(-\gamma\|u-v\|_2^2).
\]

它让距离较近的特征具有较高相似度，允许形成非线性边界。参数 C 控制对训练错误的惩罚强度，γ 控制相似度随距离衰减的速度。类别权重设为 balanced，以补偿正常图数量更多的情况。输出 G(x) 是分类决策分数，不是异常概率。

\textbf{局部分支使用正常局部记忆库。}从当前拟合部分的正常图像中随机抽取 8192 个局部向量。对待测图的每个位置 p，计算它与记忆库最相似的正常向量之间的距离：
\[
d_p(x)=\min_{m\in\mathcal M}\|z_p(x)-m\|_2.
\]

距离大，表示该位置与已见正常局部模式差异较大。代码不直接平均整张图的全部距离，而是取最大的 q 比例求平均。最终 q=0.01，即 2304 个位置中最大的 24 个距离。这是为了避免小缺陷被大量正常背景稀释，同时比只取单个最大值更能减轻孤立极值的影响。

依据：\path{Problem 2/detector.py:65}、同文件第71行、同文件第79行、同文件第93行。

\subsection{融合、选参和最终测试}

两个分支的分数尺度不同，不能直接认为“1 分”代表相同异常程度。代码用正常训练图的折外分数计算中位数与第 10—90 百分位跨度，分别校准：
\[
\widetilde B=(B-Q_{50}^B)/(Q_{90}^B-Q_{10}^B),\qquad
S=(1-\alpha)\widetilde G+\alpha\widetilde P.
\]

“折外”指分层五折中，每张图由其余四折拟合的模型评分。这样，正常图不会因为自己进入记忆库而获得不合理地小的距离，标准化统计量也不使用留出折。随后分阶段比较 4 组 SVM 参数、3 个局部聚合比例、5 个融合权重；各阶段优先比较 F1，再比较平衡准确率和 AUROC。阈值也根据训练折外分数选取。这里每类共评估 12 组分阶段候选，并非 4×3×5 的联合穷举。

最终两类都采用 C=10、γ=0.1/7168、局部权重 α=0.125。选择完成后，用全部训练图重拟合，保留折外校准参数和阈值。程序先保存冻结配置，再保存测试分数和预测，最后读官方测试标签计算指标。

依据：\path{Problem 2/detector.py:98}、同文件第115行；\path{Problem 2/run.py:66}、同文件第82行。

\subsection{怎样理解已有结果}

报告中全局与局部分支都判断正确 100/102 张，但误报图片不同；融合后正确 102/102 张。这支持“两个分支在这批图片上能够互补”，不能推导出融合在任何数据上都优于单分支。尤其全局分支 AUROC 已为 1，仍有误报：这表示异常与正常的相对排序正确，但固定阈值没有落在最佳间隙中。

报告明确说明项目早期曾用课程测试集比较版本，所以当前代码先冻结再测试，不等于整个开发过程从未接触测试结果。每类测试异常只有 15 张，新增一次漏检就会使该类召回率下降约 6.67 个百分点。此外，热图只是局部距离的可视化，没有像素级真值，不能把热图直接当作经过评价的缺陷分割结果。

依据：\path{Problem 2/course-report/sec/task2-experiments.tex:105}、同文件第140行、同文件第153行、同文件第191行、同文件第196行。\section{第三题：多变量时间序列预测}

第三题的核心是：根据此前一段时间的多个观测变量，预测下一时刻的目标变量。现有报告把 \path{OT} 定义为 CO₂ 浓度，单位为 ppm；部分代码注释仍写“室外温度”，需要统一。以下数值来自已有结果文件，本次没有重新训练。

\subsection{一条训练样本是什么}

数据有 26,200 行、21 个数值变量。程序用连续 12 行预测下一行的 \path{OT}：
\[
X_t=[x_{t-12},\ldots,x_{t-1}]\in\mathbb R^{12\times21},\qquad y_t=OT_t.
\]

例如，第 1—12 行组成第一个输入，第 13 行的 \path{OT} 是标签；窗口前移一行，就得到下一条样本。历史 \path{OT} 也在输入里，但预测时刻的 \path{OT} 不在输入中。全部数据产生 26,188 个窗口，再按索引分为训练 17,809、验证 3,142、测试 5,237 条。一个批次的输入形状是 \textnormal{(批次大小,12,21)}，输出是每个窗口对应的一个数。

若记录严格每 10 分钟一条，12 行才可解释为过去 2 小时。现有日期处理存在问题，见后文。

\subsection{为什么先标准化}

气压、温度、风速和浓度的单位、尺度不同。程序只用训练窗口估计各特征的均值与标准差，再对三部分统一变换：
\[
\widetilde x_j=(x_j-\mu_j)/s_j.
\]

目标也使用训练输入中 \path{OT} 对应的统计量缩放。模型预测后以 \(\widehat y=\mu_y+s_y\widehat{\widetilde y}\) 还原为 ppm，最后计算误差。代码将所有训练窗口展平后拟合标准化器，因此重叠窗口中的观测会重复计入统计；验证和测试数据没有参与拟合。

\subsection{两条预测分支怎样工作}

\textbf{Ridge（岭回归）}把一个 \path{12×21} 窗口展成 252 维向量，为不同变量、不同滞后位置学习线性系数，并对系数平方和施加惩罚：
\[
\min_{\beta,b}\sum_i(\widetilde y_i-\beta^Tq_i-b)^2+\alpha\|\beta\|_2^2,
\quad\alpha=1.
\]

它计算快，适合检验简单线性关系已经能解释多少变化，也能缓解相邻时间输入高度相关带来的不稳定。

\textbf{GRU}逐步读取窗口，用门控机制决定保留多少历史状态、接纳多少新信息。代码采用两层 GRU，每层隐藏维度 128；最后一步的隐藏状态经过 \path{128→64→1} 回归头输出预测。训练使用标准化目标上的均方误差、Adam、梯度裁剪和验证集早停，最多 200 轮。五个成员分别训练，预测取平均，目的是降低单个模型训练波动；是否改善仍需实验检验。LSTM 是另一个循环网络对照，已有指标文件保留了其结果。

\textbf{融合}把 GRU 集成预测 \(g\) 与 Ridge 预测 \(r\) 组合为 \(wg+(1-w)r\)。程序在验证集上搜索 0、0.05、……、1 共 21 个权重，以验证 MSE 最小者为准，再固定权重评估测试集。已有结果为 \(w=0.95\)，即 GRU 占 95\%、Ridge 占 5\%；测试标签不参与这一步选权。

\subsection{指标和已有结果怎样读}

RMSE 对大误差更敏感，MAE 表示平均偏差幅度，两者越低越好、单位均为 ppm；R² 衡量相对常数预测解释了多少目标波动，越高越好；MAPE 是平均相对误差，用百分比表示。它们强调不同方面，应结合阅读。

\begin{center}\small\begin{tabularx}{\linewidth}{Xrrrr}\toprule
方法 & RMSE & MAE & R² & MAPE（\%） \\
\midrule
持久化预测 & 6.2064 & \textbf{2.6295} & 0.9287 & \textbf{0.6006} \\
Ridge & 5.8251 & 2.8492 & 0.9372 & 0.6535 \\
LSTM & 6.2140 & 3.4342 & 0.9286 & 0.7822 \\
GRU 集成 & 5.4372 & 2.8057 & 0.9453 & 0.6375 \\
GRU + Ridge & \textbf{5.4047} & 2.7684 & \textbf{0.9460} & 0.6291 \\
\bottomrule\end{tabularx}\end{center}

融合比 GRU 集成的 RMSE 下降约 0.60\%，收益较小；持久化预测的 MAE、MAPE 反而最好。因此现有实验支持“融合降低了 RMSE”，不能支持“所有指标全面领先”。持久化基线的含义是用最近一次目标观测直接预测下一次观测。

\subsection{解释结果前必须处理的实现问题}

\par\noindent 1. \textbf{时间顺序。} \path{data_utils.py} 直接按日期字符串排序，非零填充格式会使 \path{0:50} 后接 \path{10:00}。应先用 \path{pd.to_datetime} 解析，再按时间排序，并核查重复与缺口。本次只读检查发现一处重复时间、一处 100 分钟间隔；需制定去重及补齐规则后重新建窗，才可称“过去 2 小时预测下一时刻”。现有结果对应当前错误顺序。\par\smallskip
\par\noindent 2. \textbf{缺失值。} 45 个异常占位值中，34 个位于目标列。当前全表双向插值可能使用未来观测。严格预测应按真实时间划分，采用仅依赖已知历史的填补方法，例如前向填充；保留缺失标记，缺失目标不要作为真实测试标签参与评分。\par\smallskip
\par\noindent 3. \textbf{持久化首点。} 代码以训练集末值预测测试首点，跳过了中间验证集。应直接取每个测试输入窗口最后一步的原始 \path{OT}，即 \path{X_test_raw[:, -1, target_idx]}。\par\smallskip
\par\noindent 4. \textbf{随机种子。} 当前先构建模型，再在训练函数里设种子，未控制初始化。应在每次模型构建之前设置对应种子，再训练并保存模型、划分和配置。\par\smallskip

这些修正会改变实验条件，现有成绩不能直接沿用。当前目录已有代码、结果、图表及报告；本次只解释并指出改法，没有修改队友代码或重跑实验。

来源：项目中的第三题数据处理、模型、主程序、保存指标及报告实验部分。\section{第四题：正常数据训练的甲状腺异常判断}

\subsection{与第二题的主要区别}

输入是每条记录的 6 个脱敏连续特征。训练集 1839 条全部正常，测试集 1933 条，其中正常 1839、患病 94。由于训练中没有患病示例，不能直接照第二题训练有监督二分类器，现有方法改为学习正常分布，把偏离正常模式的记录评分为异常。

代码选择第 2、3、4、6 特征，组成四维输入，按正常训练集均值与标准差标准化。测试数据复用这些统计量。特征选择来自本课程测试集分布分析，而不是已知医学机理；脱敏字段也不足以支持给它们附会具体临床含义。

依据：\path{Problem 4/main.py:22}、同文件第25行、同文件第40行；\path{Problem 4/course-report/sec/task4-experiments.tex:5}、同文件第24行、同文件第68行。

\subsection{三个模型分别观察什么}

\textbf{Isolation Forest 观察“容易被随机切分单独分出来吗”。}它反复随机选特征、切分样本。处于正常密集区域之外的点往往更容易被隔离，因此平均隔离路径较短时异常分数较高。当前设置使用 1000 棵树、每棵树使用一半输入特征，随机种子为 42。

\textbf{HBOS 观察“每个特征单独看是否罕见”。}它给各维建立直方图，估计当前取值所在区间的密度，再累加负对数密度。某维落到低密度区间，就增加异常分数。其优点是简单；局限是难以单独捕捉“每一维都常见，但组合不正常”的情况。

\textbf{MCD 观察“考虑相关关系后，离正常中心多远”。}它用稳健方式估计正常中心 μ 和协方差 Σ，再计算平方马氏距离：
\[
s_M(x)=(x-\mu)^T\Sigma^{\dagger}(x-\mu).
\]

例如两个指标正常时经常一起升高，二者同时适度升高未必异常；一个升高、另一个降低可能更不符合正常相关结构。协方差就是把这种关系纳入距离。这里的稳健估计降低部分极端正常记录对中心和散布估计的影响，但不意味着能够识别所有异常形状。

依据：\path{Problem 4/main.py:45}；\path{Problem 4/course-report/sec/task4-method.tex:46}、同文件第72行、同文件第102行。

\subsection{为什么融合名次，而不是原始分数}

三个模型分数量纲不同。代码先在当前待测批次中把各模型分数转成升序平均秩 r，再映射到 [0,1]：
\[
q_{ik}=(r_{ik}-1)/(M-1),\qquad S_i=(q_{i1}+q_{i2}+q_{i3})/3.
\]

这里 M 是批次样本数；同分样本使用平均名次。若一个样本在三个模型中都靠近异常端，融合分数就较高。秩融合减少尺度差异，但只表达本批次中的相对异常程度，不能当作患病概率。同一记录换一批对照样本，秩分数也可能变化；当前实现也不能直接用仅一个样本的批次做这项归一化。

模型中的 contamination=0.05 用于各检测器内部阈值设置；本代码最终使用连续分数再融合，并在测试标签上寻找最大 F1。因此，不能解释为“最终必定把 5\% 的人判为患病”。

依据：\path{Problem 4/main.py:33}、同文件第45行、同文件第57行、同文件第62行。

\subsection{结果和评价边界}

已有分数复算得到 AUROC=0.9846、AP=0.9060。在测试集最佳 F1 阈值约 0.948844 下，TP=81、FP=12、FN=13、TN=1827，故精确率=81/93≈0.8710、召回率=81/94≈0.8617、F1≈0.8663。代码所称 AUPRC 实际由 average\_precision\_score 计算，即 AP，不是对任意绘制的 PR 折线做梯形积分。

必须区分“拟合不使用患病标签”和“整个方案完全没有使用测试标签”：现有报告明确披露测试标签参与了特征、模型组合和阈值选择，所以这些是当前课程测试集上的探索结果。最佳 F1 不是预先固定阈值在新数据上的表现。训练集去重有 1782 条，测试中另有 56 条正常记录与训练记录相同，也会影响对独立性的判断。代码虽称固定种子，但 MCD 没有显式传入 random\_state，严格复现还需核对该随机过程及库版本。

依据：\path{Problem 4/course-report/sec/task4-experiments.tex:22}、同文件第68行、同文件第91行、同文件第128行；\path{Problem 4/main.py:49}。上述指标和重复数量已根据现有 CSV 重新核算，没有重新拟合三个检测器。\section{拓展一：图像与表格共用的异常检测框架}

\subsection{统一的究竟是什么}

这题把第二题图像和第四题表格放进一个共同流程：\textbf{特征适配→建立正常参考→计算偏离→汇总为样本分数→阈值判断}。两种模态分别建参考库和阈值，不把榛子特征与甲状腺数值混进同一近邻库，也不直接比较它们的原始分数。

NormalBagDetector 接受 N×P×D 的数组：N 个样本，每个样本 P 个描述子，每个描述子 D 维。图像有很多局部描述子；表格的一条记录可视为只有一个描述子的样本。对每个描述子计算到 k 个正常参考的平均距离，再除以 √D；图像取最大的 2\% 局部分数平均，表格 P=1 时自然退化为单个距离。这里除以 √D 只能调整尺度，不保证不同模态分数具有统一概率含义。

依据：\path{Extension 1/normal_bags.py:35}、同文件第46行、同文件第86行、同文件第98行。

\subsection{图像适配分支}

图像完整缩放至 256×256，经冻结的 ResNet18 得到 1792 维全局向量与 16×16 个 64 维局部向量。局部表示来自两层特征的平滑、拼接和固定随机投影。正常参考库最多 8192 个向量，每个位置查找 5 个最近正常参考；256 个位置中最高的 6 个距离取均值。

由于图像训练集有异常标签，最终方法增加全局 SVM，并用训练正常图的折外分数中位数和四分位距校准两个分支，再计算 S=w·校准后的SVM分数+(1−w)·校准后的局部分数。榛子最终 SVM 权重 0.5，拉链为 0.95。

这里的通用选参标准与第二题不同：SVM 按五折平均 AP、随后平均 AUROC 选参；融合权重按汇总折外 AP、随后 AUROC 选择；阈值才按折外最大 F1 选择，并列时比较精确率和较大阈值。当前拉链主方案的候选范围已固定为 C=10、γ=0.0001、w=0.95，这些设置来自之前的开发比较，不能称为当前五折从多个候选中重新选出。全量重拟合后继承折外阈值，属于经验工作点。

依据：\path{Extension 1/local_features.py:34}、同文件第74行；\path{Extension 1/supervised_refinement.py:111}、同文件第123行、同文件第132行、同文件第142行；\path{Extension 1/experiment.py:42}。

\subsection{甲状腺适配分支}

先将正常训练记录去重，再分为 1425 条拟合记录、357 条正常校准记录。拟合部分选择第 2、6 特征，学习 Yeo–Johnson 变换以调整分布形状，再以中位数及第 10—90 百分位跨度做稳健缩放。将变换后坐标记为 z₂、z₆，构造
\[
u(x)=\max(z_2,0)\max(-z_6,0).
\]

这相当于专门测量“第 2 特征向上偏离，并且第 6 特征向下偏离”的联合程度。只有两个条件同时满足，乘积才为正。正常参考是拟合描述子的中位数，当前拟合值为 0；一维曼哈顿距离因此等于 u(x) 本身。它很容易解释，但对其他方向的异常可能没有反应。

校准集不参与变换参数拟合，用来确定阈值：将 n 个正常校准分数升序排列，取第 ceil((n+1)(1−α)) 个，α=0.005 时 n=357，对应第 357 个即最大校准分数；测试分数严格大于阈值才判异常。α 越小，阈值通常越高，判为异常越保守。公式具有有限样本秩校准的形式，但本项目按已见测试结果选择过特征、方向和 α，不能因此宣称整个流程对未来数据保证 0.5\% 误报率。

依据：\path{Extension 1/power_reference.py:21}、同文件第37行、同文件第43行、同文件第62行、同文件第88行；\path{Extension 1/unified.py:12}；\path{Extension 1/experiment.py:116}。

\subsection{对照、结果和可以得出的结论}

报告比较了三类方法：基线使用全局特征的标准化、PCA 近邻和投影残差；正常参考方法使用局部或表格正常近邻；适配方法再加入上述监督融合或方向联合描述子。保留 PCA 残差，是为了避免新样本在被舍弃方向上的偏离完全消失。

最终报告榛子 55/55、拉链 47/47；甲状腺 TN=1833、FP=6、FN=1、TP=93，对应 F1=0.9637、AUROC=0.9958。移除与训练重叠的 56 条正常测试记录后，F1 不变，AUROC 为 0.9956。阈值比例由 0.005 改为 0.01 时，甲状腺误报由 6 增至 22、漏报仍为 1，说明评分排序相同也会因工作点不同产生不同判断质量。

这支持正常参考流程可承接不同特征表示，但不能把适配方法的全部提升归功于某一个公共模块：特征、监督条件和阈值都发生了变化。报告明确承认课程测试结果参与过开发选择。现有代码先保存模型、配置、预测后关联测试标签，有利于核查单次运行，仍无法消除先前开发选择造成的偏差。应把“程序可复现”“本数据上结果较好”“新数据上有效”视为三个需要分别验证的命题。

依据：\path{Extension 1/unified.py:42}、同文件第68行；\path{Extension 1/experiment.py:166}、同文件第201行；\path{Extension 1/course-report/course-report.tex:66}、同文件第68行、同文件第70行、同文件第72行、同文件第74行。
\section{拓展二：财务问答Agent的可行路线（未实施）}
这个拓展与拓展一二选一，现有组内交付选择拓展一。本节解释原题，不表示已完成财务Agent。题目提供SQLite财务数据库、9份年报HTML及56个问题，训练37题来自4家公司，测试19题来自另2家公司。

一个可复用系统应先识别问题需要“结构化数字”“年报文字”还是“组合计算”，再选择SQL查询、段落检索或受控计算工具。数值答案记录公司、财年、单位及报表口径；文字答案附支持段落。例如收入趋势问题先取连续财年收入，再计算增长率$(R_t-R_{t-1})/R_{t-1}$，最后结合年报解释变化原因。不同币种、千元／百万元、财年区间和GAAP／非GAAP口径必须核对。

训练阶段可以使用参考答案和评测反馈改进索引、工具和规则；测试前固定系统，测试参考答案与评分反馈不能反过来指导设计。最终应保存每题检索依据、SQL、计算输入输出和答案，再按包内评测脚本检验支持率。题目明确自动支持率不直接换算作业分数。最有意义的改进是把证据链和可复用流程做可靠，而非只增加工具数量。

\section{把所有题目连起来理解}
\subsection{相同之处：表示、距离、统计与验证}
第一题和第二题都从图像表示出发，但第一题寻找多个类别群组，第二题判断图像是否偏离正常或已知异常边界。第四题只有正常训练数据，因此只能先学习正常规律；拓展一抽取其中“相对正常参考的偏离”作为共用流程。第三题的输出是连续数值，重点在信息的时间可用性，不能套用随机打乱的分类评估。

距离也有不同含义。第一题的距离用于判断两张图是否相似；第二题局部分支的距离用于寻找不熟悉的局部外观；第四题马氏距离考虑特征相关性；拓展一把距离作为共同接口，但不同模态的原始数值没有天然可比性。

\subsection{监督信息和数据边界必须分层说明}
区分网络预训练、模型拟合、超参数选择、阈值确定和最终评价五个环节。某个模型拟合只用正常样本，不代表特征或阈值没有根据测试标签选择。训练内交叉验证有助于选择参数，但不能消除之前查看测试结果造成的选择偏差。时间序列的插值即使没有直接用标签训练，也可能把未来信息带入过去。

\subsection{指标并不互相替代}
异常检测中，TP是检出的异常，FP是误报的正常，FN是漏掉的异常，TN是正确识别的正常。
\[
\mathrm{Precision}=\frac{TP}{TP+FP},\quad
\mathrm{Recall}=\frac{TP}{TP+FN},\quad
F_1=\frac{2TP}{2TP+FP+FN}.
\]
精确率回答“报出的异常有多少是真的”，召回率回答“真实异常找回了多少”。AUROC综合不同阈值下的真阳性率和假阳性率，也可理解为随机异常样本的分数高于随机正常样本的概率（同分记半）。AP按召回率增量加权精确率，通常比单独准确率更适合观察异常少、正常多的样本。AUROC和AP不固定单一工作阈值，而F1依赖实际阈值。它们都需要真实标签来评价。

聚类关注ACC、ARI、NMI；异常检测同时关注排序指标AUROC/AP和阈值指标Precision/Recall/F1；回归关注RMSE/MAE等误差。全局排序完美仍可因阈值不合适出现误报，较低RMSE也不意味着较低MAE。表述结论时应指明数据、指标和条件。

\subsection{各题最值得准备的答辩问题}
\begin{enumerate}
\item 第一题：为何选这种特征？聚类类别编号怎么与真值对应？PCA是否总是有益？
\item 第二题：全局与局部分支为何互补？折外校准避免什么问题？热图是否等于分割结果？
\item 第三题：一条输入的维度是什么？未来信息如何避免泄漏？复杂模型为何必须和持久化基线比较？
\item 第四题：没有患病训练样本如何建模？马氏距离与普通欧氏距离有什么不同？阈值为什么不能在测试集上反复选择后再称盲测？
\item 拓展一：公共框架到底统一什么？哪些部分仍然依赖模态？高分来自公共模块还是特征和监督条件同时改变？
\end{enumerate}

\section*{材料与查证位置}
\addcontentsline{toc}{section}{材料与查证位置}
题目依据为 DM\_2026\_Dataset/作业说明\_2026.pdf。第一题新成果位于\path{Problem 1/}。第二、三、四题及拓展一的讲解依据对应目录的README、代码、\path{course-report/}和已保存结果。文中相对路径均以课程作业文件夹为根目录。第一题正式报告列出方法参考文献；本讲解中的其他题结果保留原报告的实验限制，不视为本轮重新训练所得。

\end{document}
